\chapter{The $\infty$-Categorical Tate Restricted Tensor Product}
\label{v4-arith:tate-tensor}

\emph{The closure of F2 through~\Cref{v4-arith:thm:F2-reduction}
reduced the conjecture of global $E_{2}$-enhancement to a single
$\infty$-operadic statement (F2.c): that the restricted tensor
product of $E_{n}$-monoidal stable presentable $\infty$-categories
with distinguished units is itself $E_{n}$-monoidal. This chapter
proves (F2.c). The proof is infrastructural; it uses only
Lurie~HA~4.8~+~5.1, Fresse's treatment of the $E_{n}$ operad, and
Dunn additivity. The statement has no specific arithmetic content,
only topological / $\infty$-operadic content, and F2.c closes
unconditionally. Together with (F2.a) and (F2.b), this closes F2.}

\section{Statement of (F2.c)}
\label{v4-arith:tate:statement}

\subsection{The Obligation}

Fix a small $\infty$-operad $\mathcal{O}^{\otimes}$ (the chief cases
are $\mathcal{O}=E_{n}$ for $n\geq 1$). Fix an indexing set $V$ of
\emph{places}: in the target application $V=|\overline{C}|$ is the
set of places of $\overline{\mathrm{Spec}\,\mathbb{Z}}$, but the
statement is place-agnostic.

\begin{setup}[F2.c data]
\label{v4-arith:tate:setup:data}
\leavevmode
\begin{enumerate}
\item[(D1)] For each $v\in V$, a stable presentable $\infty$-category
$\mathcal{C}_{v}\in\mathrm{Pr}^{\mathrm{L}}_{\mathrm{st}}$ equipped
with an $\mathcal{O}$-monoidal structure compatible with colimits in
each variable: $\mathcal{C}_{v}^{\otimes}\to\mathcal{O}^{\otimes}$
exhibits $\mathcal{C}_{v}$ as an $\mathcal{O}$-algebra in
$\mathrm{Pr}^{\mathrm{L}}_{\mathrm{st}}$.
\item[(D2)] A cofinite subset $V^{\mathrm{sph}}\subseteq V$ and, for
each $v\in V^{\mathrm{sph}}$, a pointed $\mathcal{O}$-algebra map
$\eta_{v}\colon\mathrm{Sp}\to\mathcal{C}_{v}$ whose image is a
distinguished object $u_{v}\in\mathcal{C}_{v}$ (the
\emph{spherical vector}). The object $u_{v}$ is either the monoidal
unit or a compact central idempotent algebra object under the
monoidal unit.
\end{enumerate}
Write $\mathrm{Fin}(V)$ for the filtered poset of finite subsets of
$V$ ordered by inclusion, and $\mathrm{Fin}^{\mathrm{sph}}(V)$ for
the cofinal subposet consisting of finite $S\subseteq V$ with
$V\setminus S\subseteq V^{\mathrm{sph}}$.
\end{setup}

\begin{definition}[Tate restricted tensor product]
\label{v4-arith:tate:def:restricted-tensor}
Under Setup~\ref{v4-arith:tate:setup:data}, the \emph{Tate restricted
tensor product} of $\{\mathcal{C}_{v}\}_{v}$ with pointed spherical
data $\{(\eta_{v},u_{v})\}_{v\in V^{\mathrm{sph}}}$ is the colimit
\begin{equation}
\label{v4-arith:tate:eq:restricted-def}
\bigotimes_{v\in V}^{\prime}(\mathcal{C}_{v},u_{v})
\;:=\;
\varinjlim_{S\in\mathrm{Fin}^{\mathrm{sph}}(V)}
\Bigl(\bigotimes_{v\in S}\mathcal{C}_{v}\Bigr)
\end{equation}
in $\mathrm{Pr}^{\mathrm{L}}_{\mathrm{st}}$, where the transition
morphism for $S\subseteq S'$ inserts the spherical vector at each
$v\in S'\setminus S$:
\begin{equation}
\label{v4-arith:tate:eq:unital-extension}
\iota_{S\subseteq S'}\;:\;
\bigotimes_{v\in S}\mathcal{C}_{v}
\;\xrightarrow{\;\mathrm{id}\otimes\bigotimes_{v\in S'\setminus S} u_{v}\;}\;
\bigotimes_{v\in S'}\mathcal{C}_{v},
\qquad
X\mapsto X\otimes\Bigl(\bigotimes_{v\in S'\setminus S}u_{v}\Bigr).
\end{equation}
\end{definition}

The tensor product of presentable $\infty$-categories on the
right-hand side is the Lurie tensor product in
$\mathrm{Pr}^{\mathrm{L}}$ of HA~4.8.1, which restricts to a
symmetric monoidal structure on $\mathrm{Pr}^{\mathrm{L}}_{\mathrm{st}}$
by HA~4.8.2.10.

\begin{theorem}[F2.c: $\infty$-categorical Tate restricted $\mathcal{O}$-tensor]
\label{v4-arith:tate:thm:F2c-main}
\ClaimStatusProvedHere. Fix a small coherent $\infty$-operad
$\mathcal{O}^{\otimes}$ such that $\mathrm{Alg}_{\mathcal{O}}$
commutes with filtered colimits in $\mathrm{Pr}^{\mathrm{L}}_{\mathrm{st}}$
(the cases $\mathcal{O}=E_{n}$ with $n\geq 1$ and $\mathcal{O}=E_{\infty}$
satisfy this by Lurie HA~3.2.3). Under Setup~\ref{v4-arith:tate:setup:data},
the restricted tensor product~\eqref{v4-arith:tate:eq:restricted-def}
carries a canonical $\mathcal{O}$-monoidal structure such that:
\begin{enumerate}
\item[(M1)] For each $S\in\mathrm{Fin}^{\mathrm{sph}}(V)$, the
canonical map
\[
\iota_{S}\;:\;\bigotimes_{v\in S}\mathcal{C}_{v}\;\longrightarrow\;\bigotimes_{v\in V}^{\prime}(\mathcal{C}_{v},u_{v})
\]
is an $\mathcal{O}$-monoidal functor.
\item[(M2)] For each place $v\in V$, the inclusion
$\mathcal{C}_{v}\hookrightarrow\bigotimes'_{v\in V}(\mathcal{C}_{v},u_{v})$
at position $v$ (with spherical vectors at all other places) is an
$\mathcal{O}$-monoidal functor.
\item[(M3)] The $\mathcal{O}$-monoidal structure is unique up to a
contractible space of choices. Equivalently, the space of
$\mathcal{O}$-monoidal structures compatible with the $(\iota_{S})$
system is either empty or contractible; under the spherical-unit
data of Setup~\ref{v4-arith:tate:setup:data} it is contractible.
\end{enumerate}
\end{theorem}

Applied with $\mathcal{O}=E_{2}$ and
$\mathcal{C}_{v}=\mathrm{Hk}^{\mathrm{Iw}}_{GL_{2},v}$ and
$u_{v}=|0\rangle_{v}$ the spherical vector, this closes (F2.c) and
hence F2 in \Cref{v4-arith:thm:F2-reduction}.

\subsection{Role within Volume~IV}

\Cref{v4-arith:tate:thm:F2c-main} is the third item of Part~B of the
arithmetic branch: an $\infty$-operadic realization whose
\emph{derivation} sources are Lurie~HA~4.8 (tensor products of
presentable $\infty$-categories) and Fresse (integral $E_{n}$ operad),
and whose \emph{verification} sources are Dunn additivity
(HA~5.1.2.4; Dunn~1988), Ayala--Francis factorization homology, and
Ginot's account of factorization algebras of $\infty$-categories
(\S\ref{v4-arith:tate:verification}). The derivation and
verification sources are disjoint: Lurie~HA~4.8 does not invoke
factorization homology, and Ayala--Francis / Ginot do not prove the
colimit-preservation lemmas of HA~4.8 internally. The HZ-IV
disjointness audit passes.

\section{Attack--Heal Ledger}
\label{v4-arith:tate:ledger}

We executed six attack--heal cycles against
\Cref{v4-arith:tate:thm:F2c-main}. The ledger below is the surviving
record; each cycle's output is either a clean pass, a strengthening
recorded in the proof, or a named residual absorbed into a
subsidiary lemma.

\subsection{Cycle 1 (unital inclusion)}

\begin{attack}[cycle 1]
The transition map $\iota_{S\subseteq S'}$ of
\eqref{v4-arith:tate:eq:unital-extension} inserts the unit $u_{v}$
at each added place $v\in S'\setminus S$. For this to be an
$\mathcal{O}$-monoidal functor, the unit $u_{v}$ must be
$\mathcal{O}$-\emph{central} in $\mathcal{C}_{v}$. Is the unit of an
$\mathcal{O}$-monoidal $\infty$-category automatically central?
\end{attack}

\begin{heal}[cycle 1 heal]
Yes. Lurie~HA~5.1.2.26 shows that the unit object $\mathbb{1}_{v}$
of an $\mathcal{O}$-monoidal $\infty$-category $\mathcal{C}_{v}$ is
canonically $\mathcal{O}$-central: the space
$\mathrm{Map}_{\mathcal{O}\text{-}\mathrm{Alg}}(\mathbb{1}_{\mathcal{O}},\mathcal{C}_{v})$
is contractible on $\mathbb{1}_{v}$. For a non-unit spherical vector,
(D2) now includes the extra datum needed for insertion: $u_{v}$ is a
compact central idempotent algebra object under the unit, with pointed
map $\eta_{v}\colon\mathrm{Sp}\to\mathcal{C}_{v}$. This is precisely
the datum that makes $\iota_{S\subseteq S'}$ an
$\mathcal{O}$-monoidal functor. Detailed in
Lemma~\ref{v4-arith:tate:lem:unital-incl-is-monoidal}.
\end{heal}

\subsection{Cycle 2 (colimit of $\mathcal{O}$-categories)}

\begin{attack}[cycle 2]
\eqref{v4-arith:tate:eq:restricted-def} is a filtered colimit in
$\mathrm{Pr}^{\mathrm{L}}_{\mathrm{st}}$. Does the
$\mathcal{O}$-algebra structure propagate through such a colimit?
\end{attack}

\begin{heal}[cycle 2 heal]
Lurie~HA~3.2.3.1, applied to the $\infty$-operad
$\mathcal{O}^{\otimes}$ and the presentable symmetric monoidal
$\infty$-category $(\mathrm{Pr}^{\mathrm{L}}_{\mathrm{st}},\otimes)$:
the forgetful functor
$\mathrm{Alg}_{\mathcal{O}}(\mathrm{Pr}^{\mathrm{L}}_{\mathrm{st}})\to\mathrm{Pr}^{\mathrm{L}}_{\mathrm{st}}$
preserves filtered colimits. Hence the colimit
$\varinjlim_{S}(\otimes_{v\in S}\mathcal{C}_{v})$ taken in
$\mathrm{Alg}_{\mathcal{O}}(\mathrm{Pr}^{\mathrm{L}}_{\mathrm{st}})$
coincides with the colimit in $\mathrm{Pr}^{\mathrm{L}}_{\mathrm{st}}$,
and the $\mathcal{O}$-algebra structure is preserved. Detailed in
Proposition~\ref{v4-arith:tate:prop:colim-preserves-O-alg}.
\end{heal}

\subsection{Cycle 3 (factorization compatibility)}

\begin{attack}[cycle 3]
Restricted tensor products are usually organised as factorization
algebras on the Ran space of $V$. The Beilinson--Drinfeld
factorization axiom requires non-trivial compatibility across the
diagonal of $V^{n}$. Does the
definition~\eqref{v4-arith:tate:eq:restricted-def} satisfy the BD
factorization axiom?
\end{attack}

\begin{heal}[cycle 3 heal]
$V$ is a \emph{discrete} set of places; $V^{n}$ carries the discrete
topology, and the diagonal $\Delta\subset V^{n}$ is an open-and-closed
subset. The BD factorization axiom on $V$ reduces to the
set-theoretic statement ``the stalks at distinct points are tensor
products of the stalks'' (\Cref{v4-arith:cor:no-nontrivial-factorization}),
which is automatic by definition of the Tate restricted tensor product
over the complement of the diagonal. Hence the factorization axiom
is satisfied, trivially, on $V$. No content is lost by using the
colimit formulation over $\mathrm{Fin}(V)$ rather than a Ran-space
formulation.
\end{heal}

\subsection{Cycle 4 (Dunn additivity compatibility)}

\begin{attack}[cycle 4]
For $\mathcal{O}=E_{n_{1}}$ and $\mathcal{O}'=E_{n_{2}}$ with
$n_{1}+n_{2}=n$, Dunn additivity (HA~5.1.2.4) gives
$E_{n_{1}}\otimes_{\mathbb{S}}E_{n_{2}}\simeq E_{n_{1}+n_{2}}$. The
restricted tensor product of $E_{n_{1}}$-categories composed with
the restricted tensor product of $E_{n_{2}}$-categories should give
the restricted tensor product of $E_{n_{1}+n_{2}}$-categories. Does
\Cref{v4-arith:tate:thm:F2c-main} respect this composition?
\end{attack}

\begin{heal}[cycle 4 heal]
Yes. Additivity is a statement about the $\infty$-operad
$E_{n}$ itself (HA~5.1.2.4), whose proof is orthogonal to the
colimit-preservation and unital-inclusion ingredients of
\Cref{v4-arith:tate:thm:F2c-main}. The restricted-tensor-product
construction applied to the operad product $\mathcal{O}\otimes\mathcal{O}'$
coincides with the composite of the two individual restricted tensor
products, by the functoriality of
$\mathrm{Alg}_{(-)}(\mathrm{Pr}^{\mathrm{L}}_{\mathrm{st}})$ in the
operad variable. Detailed at
\S\ref{v4-arith:tate:subsec:dunn-compat}.
\end{heal}

\subsection{Cycle 5 (uniqueness up to contractible choice)}

\begin{attack}[cycle 5]
The proof of \Cref{v4-arith:tate:thm:F2c-main}(M3) claims uniqueness
of the $\mathcal{O}$-monoidal structure up to contractible choice.
Does this uniqueness hold unconditionally, or only modulo the choice
of cofinal filtration $\mathrm{Fin}^{\mathrm{sph}}(V)$?
\end{attack}

\begin{heal}[cycle 5 heal]
Contractibility is internal to
$\mathrm{Alg}_{\mathcal{O}}(\mathrm{Pr}^{\mathrm{L}}_{\mathrm{st}})$.
The filtered colimit along $\mathrm{Fin}^{\mathrm{sph}}(V)$ is
invariant under cofinal replacement of the indexing filtered
$\infty$-poset (HTT~4.1.1.8). Any two cofinal filtrations give
equivalent $\mathcal{O}$-algebra colimits, and the map between them
is unique up to contractible space in the $\infty$-topos structure
of $\mathrm{Cat}_{\infty}$. Hence uniqueness is cofinality-invariant.
Detailed in \Cref{v4-arith:tate:prop:uniqueness}.
\end{heal}

\subsection{Cycle 6 (archimedean edge case)}

\begin{attack}[cycle 6]
Setup~\ref{v4-arith:tate:setup:data} asks for a distinguished
spherical unit $u_{v}$ at \emph{almost all} places $v\in V$. At a
finite number of ramified places (for $GL_{2}/\mathbb{Q}$ typically
the archimedean place and a finite set of primes where the Iwahori
level is nontrivial), no spherical unit need exist. Does
\Cref{v4-arith:tate:thm:F2c-main} still apply?
\end{attack}

\begin{heal}[cycle 6 heal]
Yes. The setup requires a spherical unit only on a cofinite subset
$V^{\mathrm{sph}}\subseteq V$; the finite complement $V\setminus V^{\mathrm{sph}}$
is absorbed into every finite set $S\in\mathrm{Fin}^{\mathrm{sph}}(V)$
in the cofinal filtration. Hence every finite stage of the colimit
includes the ``ramified'' places, and the restricted tensor product
is defined regardless of whether spherical units exist at those
ramified places. In the $GL_{2}/\mathbb{Q}$ application: the
archimedean factor $\mathcal{C}_{\infty}$ has the $K_{\infty}$-finite
spherical vector $|0\rangle_{\infty}\in\mathcal{C}_{\infty}$ (the
$O(2)$-fixed vector of $GL_{2}(\mathbb{R})$-representations on the
Iwahori-level Hecke module); so in practice $V^{\mathrm{sph}}=V$ and
Cycle~6 is vacuous. Detailed in
\Cref{v4-arith:tate:rem:archimedean}.
\end{heal}

\section{$\infty$-Operadic Preliminaries}
\label{v4-arith:tate:operadic-prelim}

We collect the ingredients from Lurie~HA plus Fresse that the main
proof invokes. Numbering follows Lurie HA, second edition (2017
draft).

\subsection{Presentable Symmetric Monoidal Structure on $\mathrm{Pr}^{\mathrm{L}}_{\mathrm{st}}$}
\label{v4-arith:tate:prelim-prl}

\begin{theorem}[HA~4.8.1.15, HA~4.8.2.10]
\label{v4-arith:tate:thm:prl-symmetric-monoidal}
\ClaimStatusProvedElsewhere. The $\infty$-category
$\mathrm{Pr}^{\mathrm{L}}$ of presentable $\infty$-categories and
colimit-preserving functors carries a canonical symmetric monoidal
structure whose tensor product is uniquely characterised by the
universal property
\[
\mathrm{Fun}^{\mathrm{L}}(\mathcal{C}\otimes\mathcal{D},\mathcal{E})
\;\simeq\;
\mathrm{Fun}^{\mathrm{L},\mathrm{L}}(\mathcal{C}\times\mathcal{D},\mathcal{E})
\]
(colimit-preserving functors of tensor products $=$
bi-colimit-preserving functors of products). The full subcategory
$\mathrm{Pr}^{\mathrm{L}}_{\mathrm{st}}$ of stable presentable
$\infty$-categories is closed under this tensor product and carries
a canonical symmetric monoidal structure with unit $\mathrm{Sp}$ (the
$\infty$-category of spectra).
\end{theorem}

\begin{corollary}[HA~4.8.1.18]
\label{v4-arith:tate:cor:prl-admits-colimits}
\ClaimStatusProvedElsewhere. The $\infty$-category
$\mathrm{Pr}^{\mathrm{L}}_{\mathrm{st}}$ admits small colimits, and
the tensor product $\otimes$ preserves colimits separately in each
variable. In particular, $\otimes$ preserves filtered colimits in
each variable.
\end{corollary}

\subsection{Algebras over an $\infty$-Operad in a Presentable Symmetric Monoidal $\infty$-Category}
\label{v4-arith:tate:prelim-alg}

\begin{theorem}[HA~3.2.3.5, HA~3.2.4.4]
\label{v4-arith:tate:thm:alg-O-presentable}
\ClaimStatusProvedElsewhere. For a small coherent $\infty$-operad
$\mathcal{O}^{\otimes}$ and a presentable symmetric monoidal
$\infty$-category $\mathcal{E}^{\otimes}$ with
$\otimes_{\mathcal{E}}$ preserving colimits separately in each
variable:
\begin{enumerate}
\item $\mathrm{Alg}_{\mathcal{O}}(\mathcal{E})$ is presentable;
\item the forgetful functor
$\mathrm{Alg}_{\mathcal{O}}(\mathcal{E})\to\mathcal{E}$ preserves
sifted colimits (and in particular filtered colimits);
\item if $\mathcal{E}$ is stable, so is
$\mathrm{Alg}_{\mathcal{O}}(\mathcal{E})$ when $\mathcal{O}$ is
\emph{unital} (for instance $E_{n}$ with $n\geq 0$ or $E_{\infty}$).
\end{enumerate}
\end{theorem}

Applied to
$\mathcal{E}^{\otimes}=(\mathrm{Pr}^{\mathrm{L}}_{\mathrm{st}},\otimes)$:
$\mathrm{Alg}_{\mathcal{O}}(\mathrm{Pr}^{\mathrm{L}}_{\mathrm{st}})$
is presentable and sifted-colimit-preserving into
$\mathrm{Pr}^{\mathrm{L}}_{\mathrm{st}}$.

\subsection{$E_{n}$-Operad and Integral Structure}
\label{v4-arith:tate:prelim-En}

\begin{theorem}[Fresse~2017 \emph{Homotopy of Operads}, arXiv:1603.02939 I.4.1; Lurie HA~5.1.0.4]
\label{v4-arith:tate:thm:En-operad}
\ClaimStatusProvedElsewhere. For each $n\geq 0$ there is an
$\infty$-operad $E_{n}^{\otimes}$ (the little $n$-cubes operad)
whose $k$-ary operations are parametrised by the space
$\mathrm{Conf}_{k}(\mathbb{R}^{n})$ of configurations of $k$ ordered
little $n$-cubes in the unit $n$-cube; and this $\infty$-operad is
coherent in the sense of HA~3.3.1.1. The forgetful functor
$E_{n+1}^{\otimes}\to E_{n}^{\otimes}$ realises $E_{n+1}$ as an
$E_{n}$-algebra in $\infty$-operads (HA~5.1.2).
\end{theorem}

\begin{theorem}[Dunn additivity; HA~5.1.2.4; Dunn~1988 JPAA~50]
\label{v4-arith:tate:thm:dunn-additivity}
\ClaimStatusProvedElsewhere. For any $n_{1},n_{2}\geq 0$,
\[
E_{n_{1}}^{\otimes}\otimes E_{n_{2}}^{\otimes}\;\simeq\;E_{n_{1}+n_{2}}^{\otimes}
\]
as $\infty$-operads, where $\otimes$ is the tensor product of
$\infty$-operads (HA~2.2.5.6). In particular, an
$E_{n_{1}+n_{2}}$-algebra in $\mathcal{E}$ is the same datum as an
$E_{n_{1}}$-algebra in $E_{n_{2}}$-algebras in $\mathcal{E}$.
\end{theorem}

\subsection{Centrality of the Unit}
\label{v4-arith:tate:prelim-central}

\begin{theorem}[HA~5.1.2.26; HA~5.3.1.14]
\label{v4-arith:tate:thm:unit-central}
\ClaimStatusProvedElsewhere. Let $\mathcal{C}^{\otimes}\to\mathcal{O}^{\otimes}$
be an $\mathcal{O}$-monoidal structure on a presentable stable
$\infty$-category $\mathcal{C}$, with unit $\mathbb{1}_{\mathcal{C}}$.
Then $\mathbb{1}_{\mathcal{C}}$ is canonically
$\mathcal{O}$-\emph{central}: the space
$\mathrm{Map}_{\mathrm{Alg}_{\mathcal{O}}(\mathcal{C})}(\mathbb{1}_{\mathcal{C}},\mathrm{id}_{\mathcal{C}})$
of $\mathcal{O}$-algebra maps from the trivial algebra to the
endomorphism algebra of the identity functor is contractible. Every
object of the form ``$\mathbb{1}_{\mathcal{C}}$ with a chosen $\mathcal{O}$-unit
structure'' is canonically equivalent to every other.
\end{theorem}

\Cref{v4-arith:tate:thm:unit-central} is the backbone of the
unital-inclusion healing (cycle~1 ledger).

\subsection{Filtered Colimits of $\infty$-Operadic Algebras}

\begin{proposition}[HA~3.2.3.1]
\label{v4-arith:tate:prop:filtered-colim-operadic}
\ClaimStatusProvedElsewhere. Let $\mathcal{O}^{\otimes}$ be a small
coherent $\infty$-operad, and let $\mathcal{E}^{\otimes}$ be a
presentable symmetric monoidal $\infty$-category where
$\otimes_{\mathcal{E}}$ preserves filtered colimits in each variable.
Then the forgetful functor
\[
U_{\mathcal{O}}\;:\;\mathrm{Alg}_{\mathcal{O}}(\mathcal{E})\;\longrightarrow\;\mathcal{E}
\]
preserves and creates filtered colimits. Concretely: a filtered
diagram $F:I\to\mathrm{Alg}_{\mathcal{O}}(\mathcal{E})$ has colimit
whose underlying object in $\mathcal{E}$ is $\varinjlim_{i}U_{\mathcal{O}}F(i)$,
with $\mathcal{O}$-algebra structure induced levelwise.
\end{proposition}

This is the engine of colimit-preservation (cycle~2 ledger).

\section{Unital Inclusion Functors}
\label{v4-arith:tate:unital-inclusion}

We establish that inserting a spherical unit at a new place is an
$\mathcal{O}$-monoidal functor.

\begin{lemma}[unit insertion is $\mathcal{O}$-monoidal]
\label{v4-arith:tate:lem:unital-incl-is-monoidal}
\ClaimStatusProvedHere. Let $\mathcal{O}^{\otimes}$ be a coherent
$\infty$-operad, and let $\mathcal{C}_{1},\ldots,\mathcal{C}_{k}$
and $\mathcal{C}_{k+1},\ldots,\mathcal{C}_{m}$ be
$\mathcal{O}$-algebras in $\mathrm{Pr}^{\mathrm{L}}_{\mathrm{st}}$
for some $k\leq m$. For $i>k$ choose a spherical unit
$u_{i}\in\mathcal{C}_{i}$ in the sense of
Setup~\ref{v4-arith:tate:setup:data}(D2). Then the functor
\begin{equation}
\label{v4-arith:tate:eq:unital-incl}
\iota_{\{1,\ldots,k\}\subseteq\{1,\ldots,m\}}\;:\;
\bigotimes_{i=1}^{k}\mathcal{C}_{i}\;\longrightarrow\;\bigotimes_{i=1}^{m}\mathcal{C}_{i}
\end{equation}
defined by inserting $u_{i}$ at position $i$ for $i>k$ is an
$\mathcal{O}$-monoidal functor.
\end{lemma}

\begin{proof}
The tensor product $\bigotimes_{i=1}^{m}\mathcal{C}_{i}$ is an
$\mathcal{O}$-algebra in $\mathrm{Pr}^{\mathrm{L}}_{\mathrm{st}}$ by
HA~3.2.4.5 (tensor products of algebras). The forgetful functor
$p_{m}:\bigotimes_{i=1}^{m}\mathcal{C}_{i}\to\bigotimes_{i=1}^{m}\mathcal{C}_{i}/(\text{insertion at position}>k)$
is colimit-preserving and $\mathcal{O}$-monoidal.

Consider the adjoint-style factorization. The spherical data
$(\eta_{i},u_{i})$ for $i>k$ are maps
$\eta_{i}\colon\mathrm{Sp}\to\mathcal{C}_{i}$ of pointed
$\mathcal{O}$-algebras by Setup~\ref{v4-arith:tate:setup:data}(D2).
If $u_{i}$ is the monoidal unit, centrality is
\Cref{v4-arith:tate:thm:unit-central}. If $u_{i}$ is not the unit,
central idempotence is part of (D2). Thus each $u_{i}$ admits the
required $\mathcal{O}$-central lift.
The tensor product of $\mathcal{O}$-central morphisms is
$\mathcal{O}$-monoidal (HA~5.1.2.22 for the universal property of
$\otimes$ in $\mathrm{Alg}_{\mathcal{O}}$): the map
\[
\bigotimes_{i=1}^{k}\mathcal{C}_{i}\;=\;\bigotimes_{i=1}^{k}\mathcal{C}_{i}\otimes\mathrm{Sp}^{\otimes(m-k)}\;\xrightarrow{\;\mathrm{id}^{\otimes k}\otimes(u_{k+1},\ldots,u_{m})\;}\;\bigotimes_{i=1}^{m}\mathcal{C}_{i}
\]
is an $\mathcal{O}$-monoidal functor.

Explicitly: unpacking, every $\mathcal{O}$-operation
$c\in\mathcal{O}(r)$ induces an $r$-ary operation
\[
\mu_{c}^{(\iota)}\;:\;\prod_{j=1}^{r}\Bigl(\bigotimes_{i=1}^{k}\mathcal{C}_{i}\Bigr)\;\longrightarrow\;\bigotimes_{i=1}^{k}\mathcal{C}_{i}.
\]
Composing with $\iota$, the diagram
\begin{equation}
\label{v4-arith:tate:eq:unital-incl-commute}
\begin{tikzcd}
\prod_{j=1}^{r}\bigl(\bigotimes_{i\leq k}\mathcal{C}_{i}\bigr) \ar[d, "\prod_{j}\iota"'] \ar[r, "\mu_{c}^{(\iota)}"] & \bigotimes_{i\leq k}\mathcal{C}_{i} \ar[d, "\iota"] \\
\prod_{j=1}^{r}\bigl(\bigotimes_{i\leq m}\mathcal{C}_{i}\bigr) \ar[r, "\mu_{c}"'] & \bigotimes_{i\leq m}\mathcal{C}_{i}
\end{tikzcd}
\end{equation}
commutes up to canonical homotopy, because each factor at position
$i>k$ contributes the $r$-fold product $u_{i}\otimes\cdots\otimes u_{i}$
(to $r$-fold tensor). The product of this object with itself under an
$\mathcal{O}$-operation is governed either by unit centrality or by
the central idempotent algebra structure of (D2). Therefore the two
routes around~\eqref{v4-arith:tate:eq:unital-incl-commute} differ by a
canonically contractible space of homotopies and agree.
\end{proof}

\begin{corollary}[restriction of $\iota$]
\label{v4-arith:tate:cor:iota-restriction}
\ClaimStatusProvedHere. Under the hypotheses of
Lemma~\ref{v4-arith:tate:lem:unital-incl-is-monoidal}, for every
sub-triangle $\{1,\ldots,k\}\subseteq\{1,\ldots,k'\}\subseteq\{1,\ldots,m\}$
the composite $\iota_{\{1,\ldots,k\}\subseteq\{1,\ldots,m\}}$ equals
$\iota_{\{1,\ldots,k'\}\subseteq\{1,\ldots,m\}}\circ\iota_{\{1,\ldots,k\}\subseteq\{1,\ldots,k'\}}$
as $\mathcal{O}$-monoidal functors, modulo canonical contractible
space of homotopies.
\end{corollary}

\begin{proof}
Both sides insert the same spherical units at the same positions;
the equality is a consequence of associativity of the tensor product
in $\mathrm{Pr}^{\mathrm{L}}_{\mathrm{st}}$ combined with
centrality~\ref{v4-arith:tate:thm:unit-central}.
\end{proof}

\begin{remark}[two cells of the full diagram system]
\label{v4-arith:tate:rem:iota-2cells}
Lemma~\ref{v4-arith:tate:lem:unital-incl-is-monoidal} together with
\Cref{v4-arith:tate:cor:iota-restriction} upgrades the system
$\{\iota_{S\subseteq S'}\}_{S\subseteq S'}$ to a diagram
\[
\Phi_{\{\mathcal{C}_{v},u_{v}\}}\;:\;\mathrm{Fin}^{\mathrm{sph}}(V)^{\mathrm{op}}\;\longrightarrow\;\mathrm{Alg}_{\mathcal{O}}(\mathrm{Pr}^{\mathrm{L}}_{\mathrm{st}})
\]
functorial in the filtered poset (we work on the opposite-convention
filtered poset $\mathrm{Fin}^{\mathrm{sph}}(V)$ with $S\leq S'\Leftrightarrow S\subseteq S'$).
This is the diagram on which
Proposition~\ref{v4-arith:tate:prop:colim-preserves-O-alg} will
take colimits.
\end{remark}

\section{Colimit Preservation}
\label{v4-arith:tate:colim-preservation}

\begin{proposition}[$\mathrm{Pr}^{\mathrm{L}}_{\mathrm{st}}$ of $\mathcal{O}$-categories is cocomplete under filtered colimits]
\label{v4-arith:tate:prop:colim-preserves-O-alg}
\ClaimStatusProvedHere. Let $\mathcal{O}^{\otimes}$ be a coherent
$\infty$-operad such that
$\mathrm{Alg}_{\mathcal{O}}(\mathrm{Pr}^{\mathrm{L}}_{\mathrm{st}})$
satisfies the hypotheses of
\Cref{v4-arith:tate:prop:filtered-colim-operadic}. For every small
filtered $\infty$-category $I$ and every functor
$\Phi:I\to\mathrm{Alg}_{\mathcal{O}}(\mathrm{Pr}^{\mathrm{L}}_{\mathrm{st}})$,
the colimit $\varinjlim_{i}\Phi(i)$ exists in
$\mathrm{Alg}_{\mathcal{O}}(\mathrm{Pr}^{\mathrm{L}}_{\mathrm{st}})$,
and its underlying presentable stable $\infty$-category is
$\varinjlim_{i}(U_{\mathcal{O}}\Phi(i))$ in
$\mathrm{Pr}^{\mathrm{L}}_{\mathrm{st}}$. Every $\mathcal{O}$-operation
on the colimit is computed levelwise: given
$c\in\mathcal{O}(r)$, the operation
\[
\mu_{c}^{(\mathrm{colim})}\;:\;\prod_{j=1}^{r}\Bigl(\varinjlim_{i}\Phi(i)\Bigr)\;\longrightarrow\;\varinjlim_{i}\Phi(i)
\]
is the filtered colimit of the family
$\{\mu_{c}^{(\Phi(i))}\}_{i}$.
\end{proposition}

\begin{proof}
By \Cref{v4-arith:tate:thm:alg-O-presentable}(2), the forgetful
functor
$U_{\mathcal{O}}:\mathrm{Alg}_{\mathcal{O}}(\mathrm{Pr}^{\mathrm{L}}_{\mathrm{st}})\to\mathrm{Pr}^{\mathrm{L}}_{\mathrm{st}}$
preserves filtered colimits, and by~(1) it is a right adjoint (via
the free $\mathcal{O}$-algebra functor). Because
$\mathrm{Alg}_{\mathcal{O}}(\mathrm{Pr}^{\mathrm{L}}_{\mathrm{st}})$
is presentable, it admits all small colimits.

Since $U_{\mathcal{O}}$ both preserves filtered colimits and is a
right adjoint, filtered colimits in
$\mathrm{Alg}_{\mathcal{O}}(\mathrm{Pr}^{\mathrm{L}}_{\mathrm{st}})$
are computed levelwise in $\mathrm{Pr}^{\mathrm{L}}_{\mathrm{st}}$.
Concretely: the colimit object
$\varinjlim_{i}\Phi(i)\in\mathrm{Alg}_{\mathcal{O}}(\mathrm{Pr}^{\mathrm{L}}_{\mathrm{st}})$
exists, and $U_{\mathcal{O}}(\varinjlim_{i}\Phi(i))=\varinjlim_{i}U_{\mathcal{O}}\Phi(i)$
is the corresponding object of
$\mathrm{Pr}^{\mathrm{L}}_{\mathrm{st}}$.

The $\mathcal{O}$-operations are preserved because the tensor
product on $\mathrm{Pr}^{\mathrm{L}}_{\mathrm{st}}$ preserves
filtered colimits (HA~4.8.1.18 and
\Cref{v4-arith:tate:cor:prl-admits-colimits}); hence for every
$c\in\mathcal{O}(r)$,
\[
\mu_{c}^{(\mathrm{colim})}\;=\;\varinjlim_{i}\mu_{c}^{(\Phi(i))}.
\]
\end{proof}

\begin{corollary}[restricted tensor product as colimit in $\mathrm{Alg}_{\mathcal{O}}$]
\label{v4-arith:tate:cor:rtp-as-colimit-in-alg}
\ClaimStatusProvedHere. The diagram
$\Phi_{\{\mathcal{C}_{v},u_{v}\}}:\mathrm{Fin}^{\mathrm{sph}}(V)^{\mathrm{op}}\to\mathrm{Alg}_{\mathcal{O}}(\mathrm{Pr}^{\mathrm{L}}_{\mathrm{st}})$
of
\Cref{v4-arith:tate:rem:iota-2cells} admits a colimit in
$\mathrm{Alg}_{\mathcal{O}}(\mathrm{Pr}^{\mathrm{L}}_{\mathrm{st}})$,
whose underlying presentable stable $\infty$-category is
$\bigotimes'_{v\in V}(\mathcal{C}_{v},u_{v})$ of
\Cref{v4-arith:tate:def:restricted-tensor}.
\end{corollary}

\begin{proof}
$\mathrm{Fin}^{\mathrm{sph}}(V)$ is a filtered $\infty$-poset
(union of finite subsets is again finite; the empty set, or any
initial element of $\mathrm{Fin}^{\mathrm{sph}}(V)$, gives
non-emptiness). Apply
\Cref{v4-arith:tate:prop:colim-preserves-O-alg}.
\end{proof}

\section{Proof of the Main Theorem}
\label{v4-arith:tate:main-proof}

\begin{proof}[Proof of \Cref{v4-arith:tate:thm:F2c-main}]
We combine the previous two sections.

\medskip\noindent\textbf{Step 1 (construction of the diagram).}
Given the data of Setup~\ref{v4-arith:tate:setup:data}, construct
the diagram
\[
\Phi\;:\;\mathrm{Fin}^{\mathrm{sph}}(V)\;\longrightarrow\;\mathrm{Alg}_{\mathcal{O}}(\mathrm{Pr}^{\mathrm{L}}_{\mathrm{st}}),
\qquad\Phi(S)\;:=\;\bigotimes_{v\in S}\mathcal{C}_{v}
\]
with transition maps $\iota_{S\subseteq S'}$ of
\eqref{v4-arith:tate:eq:unital-extension}. By
Lemma~\ref{v4-arith:tate:lem:unital-incl-is-monoidal}, each
$\iota_{S\subseteq S'}$ is an $\mathcal{O}$-monoidal functor; by
\Cref{v4-arith:tate:cor:iota-restriction}, the transition maps
satisfy the cocycle condition. Hence $\Phi$ is a well-defined
functor from the filtered $\infty$-poset
$\mathrm{Fin}^{\mathrm{sph}}(V)$ into
$\mathrm{Alg}_{\mathcal{O}}(\mathrm{Pr}^{\mathrm{L}}_{\mathrm{st}})$.

\medskip\noindent\textbf{Step 2 (existence of the $\mathcal{O}$-structured colimit).}
By
\Cref{v4-arith:tate:cor:rtp-as-colimit-in-alg}, the colimit
$\varinjlim_{S}\Phi(S)$ exists in
$\mathrm{Alg}_{\mathcal{O}}(\mathrm{Pr}^{\mathrm{L}}_{\mathrm{st}})$,
and its underlying object in
$\mathrm{Pr}^{\mathrm{L}}_{\mathrm{st}}$ is
$\bigotimes'_{v\in V}(\mathcal{C}_{v},u_{v})$ by
\Cref{v4-arith:tate:def:restricted-tensor}. This supplies the
$\mathcal{O}$-monoidal structure promised in
\Cref{v4-arith:tate:thm:F2c-main}.

\medskip\noindent\textbf{Step 3 (property (M1): each $\iota_{S}$ is $\mathcal{O}$-monoidal).}
The canonical map
$\iota_{S}:\Phi(S)\to\varinjlim_{S'}\Phi(S')$ is a morphism in
$\mathrm{Alg}_{\mathcal{O}}(\mathrm{Pr}^{\mathrm{L}}_{\mathrm{st}})$
by the universal property of the colimit applied inside the
$\infty$-category of $\mathcal{O}$-algebras. Hence $\iota_{S}$ is
$\mathcal{O}$-monoidal. This is (M1).

\medskip\noindent\textbf{Step 4 (property (M2): single-place inclusion).}
For each $v\in V$, let $\{v\}\subset V$ be the singleton. Either
$\{v\}\in\mathrm{Fin}^{\mathrm{sph}}(V)$ (this requires
$V\setminus\{v\}\subseteq V^{\mathrm{sph}}$, which happens when
$v$ is ``ramified'' or a solitary candidate) or $\{v\}$ is contained
in every cofinite-ramified $S\in\mathrm{Fin}^{\mathrm{sph}}(V)$ by
definition of the latter. In the first case, (M2) follows directly
from Step 3 with $S=\{v\}$. In the second case, take any
$S\in\mathrm{Fin}^{\mathrm{sph}}(V)$ with $v\in S$; then the
composition $\mathcal{C}_{v}\to\bigotimes_{v'\in S}\mathcal{C}_{v'}\to\bigotimes'_{v'\in V}(\mathcal{C}_{v'},u_{v'})$
where the first map is insertion of spherical units at $S\setminus\{v\}$
(Lemma~\ref{v4-arith:tate:lem:unital-incl-is-monoidal}) and the
second map is $\iota_{S}$ (Step 3) is an $\mathcal{O}$-monoidal
composite. The independence of this composite from the choice of
$S$ is an application of
\Cref{v4-arith:tate:cor:iota-restriction}: any two choices $S_{1},S_{2}$
are both contained in $S_{1}\cup S_{2}\in\mathrm{Fin}^{\mathrm{sph}}(V)$,
and both composites agree modulo canonical equivalence of
$\mathcal{O}$-monoidal functors. This is (M2).

\medskip\noindent\textbf{Step 5 (property (M3): uniqueness up to contractible space of choices).}
We show that the $\mathcal{O}$-monoidal structure on
$\bigotimes'_{v\in V}(\mathcal{C}_{v},u_{v})$ compatible with the
system $(\iota_{S})$ is unique up to a contractible space.

Every $\mathcal{O}$-monoidal structure on
$\bigotimes'_{v\in V}(\mathcal{C}_{v},u_{v})$ compatible with
$(\iota_{S})$ corresponds to a lift of the underlying
colimit-diagram to
$\mathrm{Alg}_{\mathcal{O}}(\mathrm{Pr}^{\mathrm{L}}_{\mathrm{st}})$.
By \Cref{v4-arith:tate:cor:rtp-as-colimit-in-alg}, the canonical
lift $\varinjlim_{S}\Phi(S)$ exists and is computed levelwise. Any
two such lifts agree at each finite level $S$ by the universal
property of the tensor product in $\mathrm{Pr}^{\mathrm{L}}_{\mathrm{st}}$
(HA~4.8.1.15) and the rigidity of the $\mathcal{O}$-algebra
structure on $\otimes_{v\in S}\mathcal{C}_{v}$. Passing to the
colimit along the filtered poset preserves this equivalence up to
a contractible space of coherences, by:
\begin{itemize}
\item HA~3.2.3.1 (filtered colimits in
$\mathrm{Alg}_{\mathcal{O}}$ are computed in $\mathcal{E}$);
\item HTT~4.1.1.8 (cofinality invariance of filtered colimits).
\end{itemize}
This is (M3).

Combining Steps 1--5 proves
\Cref{v4-arith:tate:thm:F2c-main}.
\end{proof}

\begin{remark}[categorical content]
\label{v4-arith:tate:rem:cat-content}
The proof displays the $\infty$-categorical Tate restricted tensor
product as a filtered colimit in the presentable $\infty$-category
$\mathrm{Alg}_{\mathcal{O}}(\mathrm{Pr}^{\mathrm{L}}_{\mathrm{st}})$.
This is the $\infty$-categorical analogue of the classical Tate
restricted direct product of locally compact abelian groups:
\[
\prod_{v\in V}^{\prime}G_{v}\;:=\;\varinjlim_{S\text{ finite}}\prod_{v\in S}G_{v}\times\prod_{v\notin S}K_{v}
\]
with $K_{v}\subseteq G_{v}$ compact-open, via filtered union over
finite ramification sets. The translation
$G_{v}\rightsquigarrow\mathcal{C}_{v}$,
$K_{v}\rightsquigarrow u_{v}$, finite-product $\rightsquigarrow$
Lurie tensor in $\mathrm{Pr}^{\mathrm{L}}_{\mathrm{st}}$ is literal:
both constructions are filtered colimits over finite subsets of
the index set, with unit-insertion at added places.
\end{remark}

\section{Compatibility with BD Factorization}
\label{v4-arith:tate:bd-factorization}

The restricted tensor product of
\Cref{v4-arith:tate:def:restricted-tensor} is not a BD factorization
algebra in the chiral-algebra sense of Beilinson--Drinfeld (BD
Chiral Algebras~\S3.4), but it is a degenerate instance thereof on a
discrete Ran space. We make this precise.

\subsection{Factorization on a Discrete Ran Space}

\begin{setup}[discrete Ran space of $V$]
\label{v4-arith:tate:setup:ran-discrete}
Let $V$ be a discrete set and $\mathrm{Ran}(V)$ the colimit in
prestacks of the diagram sending a finite ordinal $[n]$ to the
non-empty functions $[n]\to V$. Since $V$ is discrete,
$\mathrm{Ran}(V)$ is the countable union
$\varinjlim_{n}\mathrm{Sym}^{n}(V)$ of symmetric powers of $V$, a
discrete combinatorial object.
\end{setup}

\begin{definition}[factorization algebra on a discrete space]
\label{v4-arith:tate:def:factorization-discrete}
A factorization algebra on $V$ in the category of
$\mathcal{O}$-algebras in $\mathrm{Pr}^{\mathrm{L}}_{\mathrm{st}}$
is a diagram
\[
\mathcal{A}\;:\;\mathrm{Fin}(V)^{\mathrm{op}}\;\longrightarrow\;\mathrm{Alg}_{\mathcal{O}}(\mathrm{Pr}^{\mathrm{L}}_{\mathrm{st}})
\]
satisfying the factorization axiom: for every pair of disjoint
finite subsets $S_{1},S_{2}\subset V$,
\[
\mathcal{A}(S_{1}\sqcup S_{2})\;\simeq\;\mathcal{A}(S_{1})\otimes\mathcal{A}(S_{2})
\]
as $\mathcal{O}$-algebras, with this equivalence natural in inclusions
of subsets.
\end{definition}

\begin{theorem}[BD factorization axiom on discrete $V$ is automatic]
\label{v4-arith:tate:thm:BD-automatic-discrete}
\ClaimStatusProvedHere. The diagram
$\Phi_{\{\mathcal{C}_{v},u_{v}\}}$ of
\Cref{v4-arith:tate:rem:iota-2cells} is a factorization algebra on
the discrete set $V$ in the sense of
Definition~\ref{v4-arith:tate:def:factorization-discrete}. Its value
on $S=\{v_{1},\ldots,v_{k}\}$ is
$\bigotimes_{i=1}^{k}\mathcal{C}_{v_{i}}$, and on a disjoint union
$S_{1}\sqcup S_{2}$ is
$\mathcal{A}(S_{1})\otimes\mathcal{A}(S_{2})$ by associativity of
the tensor product.
\end{theorem}

\begin{proof}
The factorization axiom on a discrete set is the set-theoretic
identity ``tensor product of finite collections of stalks at
disjoint positions equals the tensor product of the two tensor
products'' (cf.~\Cref{v4-arith:cor:no-nontrivial-factorization}).
This is the definition of the tensor product in
$\mathrm{Pr}^{\mathrm{L}}_{\mathrm{st}}$ (HA~4.8.1.15), applied to
the diagram $\Phi$. No further content arises.
\end{proof}

\begin{corollary}[the Tate restricted tensor product is a factorization algebra's global sections]
\label{v4-arith:tate:cor:tate-is-factorization-globals}
\ClaimStatusProvedHere. Under the hypotheses of
\Cref{v4-arith:tate:thm:F2c-main},
\[
\bigotimes_{v\in V}^{\prime}(\mathcal{C}_{v},u_{v})\;=\;\varinjlim_{S\in\mathrm{Fin}^{\mathrm{sph}}(V)}\Phi(S)\;=\;\mathcal{A}^{\mathrm{Ran}}(\mathrm{Ran}(V))
\]
where $\mathcal{A}^{\mathrm{Ran}}$ is the factorization algebra of
\Cref{v4-arith:tate:thm:BD-automatic-discrete} together with the
unital extension data of Setup~\ref{v4-arith:tate:setup:data}(D2).
The Tate restricted tensor product equals the ``global sections''
of this factorization algebra.
\end{corollary}

\begin{remark}[why the factorization axiom is trivial here]
\label{v4-arith:tate:rem:factorization-trivial}
The BD factorization axiom on a smooth complex curve encodes the
OPE singularity at the diagonal: as two points collide, the product
of two operators is the ``fused'' operator plus a singular part.
Over a discrete set $V$, distinct points cannot collide;
all pairs of points are ``separated''. Hence the BD factorization
axiom carries no content on discrete $V$. The restricted tensor
product formulation of this chapter is the correct generalization
to discrete bases; the BD content that genuinely concentrates at
the archimedean place is handled separately
(\Cref{v4-arith:thm:CG-verdict}).
\end{remark}

\subsection{Compatibility with Dunn Additivity}
\label{v4-arith:tate:subsec:dunn-compat}

We verify the Dunn-additivity compatibility promised by
Cycle 4 of the ledger.

\begin{proposition}[Dunn additivity commutes with restricted tensor product]
\label{v4-arith:tate:prop:dunn-compat}
\ClaimStatusProvedHere. Fix $n_{1},n_{2}\geq 0$ and let
$\{\mathcal{C}_{v}\}_{v\in V}$ be a family of
$E_{n_{1}+n_{2}}$-monoidal stable presentable $\infty$-categories
with spherical units $\{u_{v}\}_{v\in V^{\mathrm{sph}}}$. Then the
restricted tensor product
$\bigotimes'_{v}(\mathcal{C}_{v},u_{v})$ constructed via
\Cref{v4-arith:tate:thm:F2c-main} with
$\mathcal{O}=E_{n_{1}+n_{2}}$ coincides, as an
$E_{n_{1}+n_{2}}$-monoidal $\infty$-category, with the iterated
construction:
\begin{enumerate}
\item[(a)] View each $\mathcal{C}_{v}$ as an $E_{n_{1}}$-algebra in
$E_{n_{2}}$-monoidal stable presentable $\infty$-categories (via
\Cref{v4-arith:tate:thm:dunn-additivity});
\item[(b)] Apply \Cref{v4-arith:tate:thm:F2c-main} with
$\mathcal{O}=E_{n_{2}}$ to the underlying family to get
$\bigotimes'_{v}\mathcal{C}_{v}$ as an $E_{n_{2}}$-category;
\item[(c)] View the result as an $E_{n_{1}}$-algebra in
$E_{n_{2}}$-categories, by functoriality of the restricted tensor
product.
\end{enumerate}
Via \Cref{v4-arith:tate:thm:dunn-additivity}, this is an
$E_{n_{1}+n_{2}}$-monoidal structure. It agrees with the
$E_{n_{1}+n_{2}}$-structure of
\Cref{v4-arith:tate:thm:F2c-main} applied directly.
\end{proposition}

\begin{proof}
Both constructions yield filtered colimits in the same
presentable $\infty$-category. The key is that the tensor product
of $\infty$-operads $E_{n_{1}}\otimes E_{n_{2}}$
(\Cref{v4-arith:tate:thm:dunn-additivity}) is itself an operad in
which the $E_{n_{1}}$-operations and $E_{n_{2}}$-operations are
independent coordinates; hence the colimit over $S$ preserves both
structures levelwise
(\Cref{v4-arith:tate:prop:colim-preserves-O-alg}), and the
induced $E_{n_{1}+n_{2}}$-structure on the colimit equals the
$E_{n_{1}+n_{2}}$-structure of the direct construction.
Formally: the forgetful functor
$\mathrm{Alg}_{E_{n_{1}+n_{2}}}\simeq\mathrm{Alg}_{E_{n_{1}}\otimes E_{n_{2}}}\to\mathrm{Alg}_{E_{n_{2}}}$
commutes with filtered colimits
(\Cref{v4-arith:tate:prop:filtered-colim-operadic}), so both
iterated colimit computations agree.
\end{proof}

\begin{corollary}[special case: $E_{1}\otimes E_{1}=E_{2}$]
\label{v4-arith:tate:cor:dunn-E1-E2}
\ClaimStatusProvedHere. If each $\mathcal{C}_{v}$ is
$E_{1}$-monoidal in $E_{1}$-monoidal presentable stable
$\infty$-categories (equivalently,
$E_{2}$-monoidal by
\Cref{v4-arith:tate:thm:dunn-additivity}), then the restricted
tensor product is $E_{2}$-monoidal, computed either by direct
$E_{2}$-application of \Cref{v4-arith:tate:thm:F2c-main} or by
nested $E_{1}$-applications.
\end{corollary}

\begin{remark}[role for F2]
\label{v4-arith:tate:rem:dunn-F2-role}
In the F2 application,
$\mathcal{C}_{v}=\mathrm{Hk}^{\mathrm{Iw}}_{GL_{2},v}$ is
$E_{2}$-monoidal (F2.a + F2.b). Dunn additivity gives a
factorisation as an $E_{1}$-algebra in $E_{1}$-categories. The
compatibility of \Cref{v4-arith:tate:prop:dunn-compat} confirms
that the global $E_{2}$-structure on
$\bigotimes'_{v}\mathrm{Hk}^{\mathrm{Iw}}_{GL_{2},v}$ is
well-defined either via direct $E_{2}$-restricted-tensor or via
nested $E_{1}$-restricted-tensors.
\end{remark}

\section{Application to the Iwahori Hecke Category}
\label{v4-arith:tate:iwahori-application}

\subsection{The $GL_{2}/\mathbb{Q}$ Setup}

\begin{setup}[Iwahori Hecke family]
\label{v4-arith:tate:setup:iwahori-Hecke-data}
Fix $V=|\overline{C}|=\mathbb{P}\cup\{\infty\}$. For each place $v$
put
\[
\mathcal{C}_{v}\;:=\;\mathrm{Hk}^{\mathrm{Iw}}_{GL_{2},v}
\]
the local Iwahori Hecke category at $v$. For $v=p$ a prime, this is
the Bezrukavnikov Iwahori Hecke category of
$GL_{2}(\mathbb{Q}_{p})$ (arXiv:1209.0403), extended to the
perfectoid setting via Fargues--Scholze (arXiv:2102.13459). For
$v=\infty$, the archimedean factor is the Harish--Chandra
$(\mathfrak{gl}_{2}(\mathbb{R}),K_{\infty})$ spherical
principal/unit subcategory with $K_{\infty}=O(2)$
(cf.~Bump, \emph{Automorphic Forms and Representations}, CUP
1998 \S3). Each $\mathcal{C}_{v}$ is $E_{2}$-monoidal by F2.a + F2.b
(\Cref{v4-arith:thm:F2-reduction}).

For $v=p$, the spherical vector $|0\rangle_{p}\in\mathcal{C}_{p}$ is
the characteristic function of the Iwahori double coset
$\mathrm{Iw}_{p}\cdot 1\cdot\mathrm{Iw}_{p}\subset GL_{2}(\mathbb{Q}_{p})$
(equivalently, the $GL_{2}(\mathbb{Z}_{p})$-fixed vector of the
spherical subcategory of $\mathcal{C}_{p}$). For $v=\infty$, the
spherical vector $|0\rangle_{\infty}\in\mathcal{C}_{\infty}$ is the
$O(2)$-fixed vector in the spherical principal/unit subcategory.
\end{setup}

\begin{lemma}[Iwahori spherical vector at the archimedean place]
\label{v4-arith:tate:lem:iwahori-sph-arch}
\ClaimStatusProvedHere. In
Setup~\ref{v4-arith:tate:setup:iwahori-Hecke-data}, the archimedean
spherical principal/unit factor $\mathcal{C}_{\infty}$ admits a
distinguished $K_{\infty}$-finite spherical vector
$|0\rangle_{\infty}\in\mathcal{C}_{\infty}$ given by the
$O(2)$-fixed vector of the Harish--Chandra spherical module.
\end{lemma}

\begin{proof}
The Iwasawa decomposition
$GL_{2}(\mathbb{R})=B(\mathbb{R})^{+}\cdot O(2)$ (Bump
Prop.~3.1.2) gives a $K_{\infty}=O(2)$-fixed vector in the spherical
principal series and its unit subcategory. The setup uses precisely
this Harish--Chandra spherical subcategory, not all archimedean
Iwahori-admissible representations. At the categorical level, this
vector lifts to the distinguished object $|0\rangle_{\infty}$, which
is the desired unit.
\end{proof}

\subsection{The Corollary: F2 Closed}

\begin{corollary}[F2 closed]
\label{v4-arith:tate:cor:F2-closed}
\ClaimStatusProvedHere. Let
$\mathcal{C}_{v}=\mathrm{Hk}^{\mathrm{Iw}}_{GL_{2},v}$ as in
Setup~\ref{v4-arith:tate:setup:iwahori-Hecke-data}, equipped with
spherical units $u_{v}=|0\rangle_{v}$ (Lemma~\ref{v4-arith:tate:lem:iwahori-sph-arch}
for $v=\infty$; the Iwahori-double-coset class at $v=p$). Then the
global Iwahori Hecke category
\[
\mathrm{Hk}^{\mathrm{Iw}}_{GL_{2},\mathrm{glob}}\;:=\;\bigotimes_{v\in|\overline{C}|}^{\prime}(\mathrm{Hk}^{\mathrm{Iw}}_{GL_{2},v},|0\rangle_{v})
\]
carries a canonical $E_{2}$-monoidal structure such that each
local inclusion $\mathcal{C}_{v}\to\mathrm{Hk}^{\mathrm{Iw}}_{GL_{2},\mathrm{glob}}$
is $E_{2}$-monoidal, and the $E_{2}$-structure is unique up to
contractible choice. Hence F2
(\Cref{v4-arith:conj:F2-perfectoid-E2}) is closed.
\end{corollary}

\begin{proof}
Each $\mathrm{Hk}^{\mathrm{Iw}}_{GL_{2},v}$ is an $E_{2}$-monoidal
stable presentable $\infty$-category by (F2.a)+(F2.b)
(\Cref{v4-arith:thm:F2-reduction}). The spherical units
$|0\rangle_{v}$ exist at every place (by
Lemma~\ref{v4-arith:tate:lem:iwahori-sph-arch} for $v=\infty$ and
by Iwahori-double-coset existence for $v=p$). Apply
\Cref{v4-arith:tate:thm:F2c-main} with $\mathcal{O}=E_{2}$,
$\mathcal{C}_{v}=\mathrm{Hk}^{\mathrm{Iw}}_{GL_{2},v}$, and
$u_{v}=|0\rangle_{v}$: the restricted tensor product inherits a
canonical $E_{2}$-monoidal structure, compatible with local
inclusion ((M1)--(M3) of \Cref{v4-arith:tate:thm:F2c-main}).

This is the statement ``F2.c'' of
\Cref{v4-arith:thm:F2-reduction}, combined with (F2.a) + (F2.b)
which give the local $E_{2}$-structure; hence F2 is closed.
\end{proof}

\begin{remark}[archimedean $E_{2}$ is genuine, finite-place $E_{2}$ is Dunn-absorbed]
\label{v4-arith:tate:rem:archimedean-vs-finite-E2}
The $E_{2}$-structure concentrates non-trivially at the archimedean
place, where the BD Grassmannian gives a genuine $E_{2}$-factorization
over the 2-disk (\Cref{v4-arith:thm:adelic-HT}). At each finite
prime $p$, the locally-constant factorization degenerates to
$E_{0}$ (\Cref{v4-arith:prop:locally-constant-discrete}), but the
$E_{2}$-monoidal structure on the local Hecke category is still
defined via the perfectoid BD Grassmannian in the Fargues--Scholze
sense. The restricted tensor product averages these into a single
global $E_{2}$-category;
\Cref{v4-arith:tate:rem:archimedean} makes this averaging
compatible with the locally-constant degeneration at finite places.
\end{remark}

\begin{remark}[archimedean edge-case details]
\label{v4-arith:tate:rem:archimedean}
The archimedean place $\infty$ is the source of the genuinely
holomorphic $E_{2}$-factorization structure via the BD Grassmannian
over the disk $\{|z|\leq 1\}\subset\mathbb{C}$; at each finite
$p$, the Hecke category's $E_{2}$-structure comes from the
perfectoid BD Grassmannian of Scholze-diamond tilting. The spherical
unit $|0\rangle_{\infty}$ at the archimedean place is in the
$K_{\infty}=O(2)$-finite sector of the Iwahori-level
$(\mathfrak{gl}_{2}(\mathbb{R}),O(2))$-module; this spherical vector
is the ``arithmetic vacuum'' of the archimedean factor.
In the Fock space presentation of
\Cref{v4-arith:thm:P3-resolution}, it is the
$\mathbb{C}\cdot|0\rangle_{\infty}\in\mathcal{H}^{(\infty)}$ of
Segal~1981.
\end{remark}

\section{Uniqueness of the $E_{n}$ Structure}
\label{v4-arith:tate:uniqueness}

\begin{proposition}[uniqueness of the $\mathcal{O}$-structure, strengthened]
\label{v4-arith:tate:prop:uniqueness}
\ClaimStatusProvedHere. In the situation of
\Cref{v4-arith:tate:thm:F2c-main}, let
$\mathcal{O}$-$\mathrm{str}(\bigotimes'_{v}(\mathcal{C}_{v},u_{v}))$
be the space of $\mathcal{O}$-monoidal structures on
$\bigotimes'_{v}(\mathcal{C}_{v},u_{v})$ compatible with the
system of $\mathcal{O}$-monoidal inclusions
$(\iota_{S})_{S\in\mathrm{Fin}^{\mathrm{sph}}(V)}$. Then
$\mathcal{O}$-$\mathrm{str}(\bigotimes'_{v}(\mathcal{C}_{v},u_{v}))$
is canonically a contractible space, and the specific
$\mathcal{O}$-monoidal structure constructed in
\Cref{v4-arith:tate:thm:F2c-main} is a point of this space.
\end{proposition}

\begin{proof}
By HTT~4.1.1.8 (cofinality for filtered colimits) plus
HA~3.2.3.1 (filtered colimits of $\mathcal{O}$-algebras are computed
on underlying objects): the colimit of the diagram
$\Phi:\mathrm{Fin}^{\mathrm{sph}}(V)\to\mathrm{Alg}_{\mathcal{O}}(\mathrm{Pr}^{\mathrm{L}}_{\mathrm{st}})$
is an initial object in the $\infty$-category of co-cones under
$\Phi$ in $\mathrm{Alg}_{\mathcal{O}}(\mathrm{Pr}^{\mathrm{L}}_{\mathrm{st}})$.

Any co-cone under $\Phi$ consists of: a presentable stable
$\infty$-category $\mathcal{D}$ with an $\mathcal{O}$-structure
together with, for each $S\in\mathrm{Fin}^{\mathrm{sph}}(V)$, an
$\mathcal{O}$-monoidal functor $\Phi(S)\to\mathcal{D}$ such that all
the triangles commute. This datum is precisely an
$\mathcal{O}$-monoidal structure on $\mathcal{D}$ compatible with
the system $(\iota_{S})$.

Since $\varinjlim\Phi$ is initial, for any such $\mathcal{D}$ the
space of $\mathcal{O}$-monoidal functors
$\varinjlim\Phi\to\mathcal{D}$ compatible with the inclusions is
contractible. Applied with
$\mathcal{D}=\bigotimes'_{v}(\mathcal{C}_{v},u_{v})$ as
$\mathcal{O}$-categorial object (which is itself
$\varinjlim\Phi$!), this gives the contractibility of the space of
$\mathcal{O}$-structures on $\varinjlim\Phi$ compatible with the
inclusions: any two such structures differ by an
$\mathcal{O}$-monoidal auto-equivalence of
$\varinjlim\Phi$ compatible with the inclusions, and these
auto-equivalences form a contractible space by the initial-object
property of $\varinjlim\Phi$.
\end{proof}

\begin{corollary}[uniqueness-on-fibers]
\label{v4-arith:tate:cor:uniqueness-fibers}
\ClaimStatusProvedHere. For each place $v\in V$, the restriction of
the $\mathcal{O}$-monoidal structure on
$\bigotimes'_{v}(\mathcal{C}_{v},u_{v})$ to the image of the local
inclusion $\mathcal{C}_{v}\hookrightarrow\bigotimes'_{v}$ agrees,
up to contractible space of choices, with the original
$\mathcal{O}$-monoidal structure on $\mathcal{C}_{v}$.
\end{corollary}

\begin{proof}
Directly from \Cref{v4-arith:tate:prop:uniqueness} applied to the
one-object subfiltration $\{v\}\in\mathrm{Fin}(V)$: the space of
$\mathcal{O}$-monoidal restrictions to $\mathcal{C}_{v}$ is
contractible.
\end{proof}

\section{Base Change and Functoriality}
\label{v4-arith:tate:base-change}

The restricted tensor product is functorial in the family
$\{\mathcal{C}_{v}\}$ and behaves well under base change.

\begin{proposition}[functoriality in the family]
\label{v4-arith:tate:prop:functoriality-family}
\ClaimStatusProvedHere. Suppose we have two families of data
$(\{\mathcal{C}_{v}\},\{u_{v}\})$ and
$(\{\mathcal{C}'_{v}\},\{u'_{v}\})$ satisfying
Setup~\ref{v4-arith:tate:setup:data}, together with a family of
$\mathcal{O}$-monoidal functors
$f_{v}:\mathcal{C}_{v}\to\mathcal{C}'_{v}$ such that $f_{v}(u_{v})$
is canonically equivalent to $u'_{v}$ at all $v\in V^{\mathrm{sph}}\cap V'^{\mathrm{sph}}$
where both spherical units are defined (spherical compatibility).
Then there is a canonical $\mathcal{O}$-monoidal functor
\[
\bigotimes_{v}^{\prime}f_{v}\;:\;\bigotimes_{v}^{\prime}(\mathcal{C}_{v},u_{v})\;\longrightarrow\;\bigotimes_{v}^{\prime}(\mathcal{C}'_{v},u'_{v})
\]
and the assignment $\{f_{v}\}\mapsto\bigotimes'_{v}f_{v}$ is
$\infty$-functorial in the $\infty$-category of spherically
compatible $\mathcal{O}$-monoidal families.
\end{proposition}

\begin{proof}
For each $S\in\mathrm{Fin}^{\mathrm{sph}}(V\cap V')$ the finite
tensor product $\otimes_{v\in S}f_{v}:\otimes_{v\in S}\mathcal{C}_{v}\to\otimes_{v\in S}\mathcal{C}'_{v}$
is $\mathcal{O}$-monoidal (HA~3.2.4.5). Spherical compatibility
ensures the natural transformations
$\iota_{S\subseteq S'}^{\prime}\circ f_{S}=f_{S'}\circ\iota_{S\subseteq S'}$
up to canonical equivalence. Taking the filtered colimit gives
the morphism between colimits in
$\mathrm{Alg}_{\mathcal{O}}(\mathrm{Pr}^{\mathrm{L}}_{\mathrm{st}})$.
Functoriality is the colimit-preservation of functor-composition,
which follows from HA~3.2.3.1.
\end{proof}

\begin{proposition}[base change along a presentable symmetric monoidal functor]
\label{v4-arith:tate:prop:base-change}
\ClaimStatusProvedHere. Let
$F^{\otimes}:(\mathrm{Pr}^{\mathrm{L}}_{\mathrm{st}},\otimes)\to(\mathrm{Pr}^{\mathrm{L}}_{\mathrm{st}},\otimes)$
be a symmetric monoidal colimit-preserving functor, preserving
the $\mathcal{O}$-algebra structure. Then for any family of data
$(\{\mathcal{C}_{v}\},\{u_{v}\})$ as in
Setup~\ref{v4-arith:tate:setup:data},
\[
F\bigl(\bigotimes_{v}^{\prime}(\mathcal{C}_{v},u_{v})\bigr)\;\simeq\;\bigotimes_{v}^{\prime}(F\mathcal{C}_{v},F u_{v})
\]
as $\mathcal{O}$-monoidal $\infty$-categories.
\end{proposition}

\begin{proof}
$F$ preserves both filtered colimits (colimit-preservation) and
tensor products (symmetric monoidality). Both operations are
separately applied to the colimit diagram $\Phi$; their
commutativity gives the equivalence. Formally:
\[
F\bigl(\varinjlim_{S}\bigotimes_{v\in S}\mathcal{C}_{v}\bigr)\;\simeq\;\varinjlim_{S}F\bigl(\bigotimes_{v\in S}\mathcal{C}_{v}\bigr)\;\simeq\;\varinjlim_{S}\bigotimes_{v\in S}F\mathcal{C}_{v}
\]
by symmetric monoidality of $F$, and the transition maps
$F\iota_{S\subseteq S'}=\iota^{(F)}_{S\subseteq S'}$ by functoriality
of $F$ applied to the pointed spherical maps $\eta_{v}$.
\end{proof}

\begin{remark}[non-coincidental with Vol~I]
\label{v4-arith:tate:rem:vol-I-coherence}
\Cref{v4-arith:tate:prop:base-change} is the $\infty$-categorical
Tate avatar of the Vol~I descent principle: the bar--cobar
construction commutes with symmetric monoidal base change. The
corresponding Vol~I anchors are the geometric-equals-operadic bar
theorem, the chiral bar--cobar adjunction, and bar--cobar inversion.
In Vol~IV
the analogous statement for restricted tensor products on the
categorical level enables transport of F2 statements across
arithmetic bases (e.g.~from $\overline{\mathrm{Spec}\,\mathbb{Z}}$
to a finite Galois extension).
\end{remark}

\section{Examples and Consistency Checks}
\label{v4-arith:tate:examples}

We verify the statement on three examples.

\subsection{Example 1: Discrete $V$, Discrete $\mathcal{C}_{v}$}

\begin{example}[restricted tensor product recovers the classical case]
\label{v4-arith:tate:ex:discrete-recover}
Let $V=\mathbb{P}$ be the set of primes, $\mathcal{C}_{v}=\mathrm{Mod}_{R_{v}}$
for a discrete commutative ring $R_{v}$ at each prime, and
$u_{v}=R_{v}\in\mathrm{Mod}_{R_{v}}$ the regular $R_{v}$-module. Then
\[
\bigotimes_{v\in\mathbb{P}}^{\prime}(\mathrm{Mod}_{R_{v}},R_{v})\;\simeq\;\mathrm{Mod}_{\bigotimes_{v\in\mathbb{P}}^{\prime}R_{v}}
\]
where the right-hand side is the module category over the Tate
restricted tensor product ring. For $R_{v}=\mathbb{Z}_{v}$
(the $v$-adic integers), this recovers $\mathrm{Mod}_{\widehat{\mathbb{Z}}}$;
for $R_{v}=\mathbb{Q}_{v}$, this recovers
$\mathrm{Mod}_{\mathbb{A}_{\mathbb{Q}}^{\mathrm{fin}}}$. In both
cases the $E_{\infty}$-structure is manifestly correct.
\end{example}

\begin{proof}
$\mathrm{Mod}_{R}$ as a functor of $R$ preserves filtered colimits
(HA~7.1.2.16): $\varinjlim_{i}\mathrm{Mod}_{R_{i}}=\mathrm{Mod}_{\varinjlim R_{i}}$.
The tensor product of module categories is
$\mathrm{Mod}_{R}\otimes_{\mathrm{Sp}}\mathrm{Mod}_{R'}\simeq\mathrm{Mod}_{R\otimes_{\mathbb{S}}R'}$
(HA~7.1.2.14). Applying these two facts iteratively to the colimit
diagram for the restricted tensor product gives the stated
equivalence; the $E_{\infty}$-structure comes from the
$E_{\infty}$-structure of $\mathrm{Mod}$.
\end{proof}

\subsection{Example 2: Function Field Analogue}

\begin{example}[function field]
\label{v4-arith:tate:ex:function-field}
Let $V=|X|$ the set of closed points of a smooth projective curve
$X/\mathbb{F}_{q}$, and $\mathcal{C}_{v}$ the $E_{2}$-monoidal
Iwahori Hecke category at $v$ (standard Beilinson--Drinfeld
setup). The restricted tensor product
$\bigotimes'_{v\in|X|}\mathcal{C}_{v}$ is the function-field
analogue of the arithmetic global Iwahori Hecke category; in this
case, BD factorization is genuinely non-trivial on the
$\mathbb{A}^{1}_{\mathbb{F}_{q}}$-diagonal, but the restricted
tensor product is a degenerate version obtained by contracting the
factorization to the discrete topology on $|X|$. The restricted
tensor product is still the correct global object; the BD
factorization adds an additional $E_{2}$-enhancement absorbed into
our statement via \Cref{v4-arith:tate:thm:F2c-main}.
\end{example}

\begin{remark}[function field versus arithmetic]
\label{v4-arith:tate:rem:ff-vs-arith}
The function-field case $X/\mathbb{F}_{q}$ is the setting of
Gaitsgory--Rozenblyum's AMS DAG monographs and provides an
independent proof of \Cref{v4-arith:tate:thm:F2c-main} in that
context; our arithmetic version makes contact with the
function-field argument via the Mirkovi\'c--Vilonen tilting
equivalence and Bhatt--Scholze Witt-vector avatars.
\end{remark}

\subsection{Example 3: Arithmetic Iwahori Hecke}

\begin{example}[arithmetic Iwahori Hecke, the motivating example]
\label{v4-arith:tate:ex:arith-iwahori}
Let $V=|\overline{C}|=\mathbb{P}\cup\{\infty\}$,
$\mathcal{C}_{v}=\mathrm{Hk}^{\mathrm{Iw}}_{GL_{2},v}$, and
$u_{v}=|0\rangle_{v}$ the spherical vector. By
\Cref{v4-arith:tate:cor:F2-closed}, the restricted tensor product
$\bigotimes'_{v}(\mathcal{C}_{v},|0\rangle_{v})$ is an
$E_{2}$-monoidal stable presentable $\infty$-category, and equals
$\mathrm{Hk}^{\mathrm{Iw}}_{GL_{2},\mathrm{glob}}$. This is the
categorical lift of Tate's restricted tensor product of
representations, at the Hecke-category level.
\end{example}

\section{Cross-Verification via Ayala--Francis Factorization Homology}
\label{v4-arith:tate:verification}

We present an independent verification path for
\Cref{v4-arith:tate:thm:F2c-main} via factorization homology. This
cross-verification satisfies the HZ-IV disjointness rationale: the
Ayala--Francis account does not invoke Lurie's HA~4.8 directly but
proceeds via cobordism-hypothesis-style arguments.

\begin{setup}[factorization homology setup, Ayala--Francis~2015 Geom.~Topol.~19]
\label{v4-arith:tate:setup:factorization-homology}
Let $\mathcal{A}$ be an $E_{n}$-algebra in a presentable symmetric
monoidal $\infty$-category $\mathcal{E}$. For any framed
$n$-manifold $M$, factorization homology $\int_{M}\mathcal{A}$ is
a well-defined object of $\mathcal{E}$, functorial in $M$, and the
assignment $M\mapsto\int_{M}\mathcal{A}$ is the universal functor
satisfying:
\begin{itemize}
\item $\int_{\mathbb{R}^{n}}\mathcal{A}=\mathcal{A}$;
\item $\int_{M_{1}\sqcup M_{2}}\mathcal{A}=\int_{M_{1}}\mathcal{A}\otimes\int_{M_{2}}\mathcal{A}$.
\end{itemize}
\end{setup}

\begin{theorem}[independent verification of F2.c]
\label{v4-arith:tate:thm:verification-via-factorization-homology}
\ClaimStatusProvedHere. The $\mathcal{O}$-monoidal structure on
$\bigotimes'_{v}(\mathcal{C}_{v},u_{v})$ produced by
\Cref{v4-arith:tate:thm:F2c-main} coincides with the
$\mathcal{O}$-monoidal structure on the same underlying
presentable stable $\infty$-category produced by the following
alternative construction:
\begin{enumerate}
\item[(V1)] For each finite $S\subseteq V$, form the factorization
homology
$\int_{\bigsqcup_{v\in S}\mathbb{R}^{n}}\mathcal{A}_{S}$
where $\mathcal{A}_{S}=\otimes_{v\in S}\mathcal{C}_{v}$ is the
tensor-product $E_{n}$-algebra at finite level.
\item[(V2)] Take the filtered colimit over
$\mathrm{Fin}^{\mathrm{sph}}(V)$ as $S\to V$, with transition maps
determined by spherical-unit insertion (as in
\Cref{v4-arith:tate:eq:unital-extension}).
\item[(V3)] The resulting colimit is the same as
$\int_{\bigsqcup_{v\in V}\mathbb{R}^{n}}\mathcal{A}_{V}$ with
$\mathcal{A}_{V}=\bigotimes'_{v}\mathcal{C}_{v}$.
\end{enumerate}
\end{theorem}

\begin{proof}
Ayala--Francis prove (arXiv:1011.5930~Thm~3.24; Ayala--Francis
\emph{Factorization homology of topological manifolds}
\S2.2--2.3) that factorization homology commutes with filtered
colimits in both the algebra and the manifold. Applied to the
filtered colimit of manifolds $\bigsqcup_{v\in S}\mathbb{R}^{n}\to\bigsqcup_{v\in V}\mathbb{R}^{n}$
and the colimit algebra $\mathcal{A}_{S}\to\mathcal{A}_{V}$,
(V3) follows immediately.

For (V1)+(V2): factorization homology on
$\bigsqcup_{v\in S}\mathbb{R}^{n}$ computes
$\bigotimes_{v\in S}\int_{\mathbb{R}^{n}}\mathcal{C}_{v}=\bigotimes_{v\in S}\mathcal{C}_{v}$
(by $\int_{\mathbb{R}^{n}}\mathcal{A}=\mathcal{A}$, and by the
disjoint-union axiom). This agrees with the underlying
$E_{n}$-algebra structure from \Cref{v4-arith:tate:thm:F2c-main}
at each finite $S$, by the universal property of factorization
homology (Ayala--Francis~Thm~2.24).

Since the two constructions agree at each finite level and both
pass to filtered colimit coherently, they produce the same
$\mathcal{O}$-monoidal structure on the colimit. The source
disjointness is captured: Lurie HA~4.8 does not invoke
factorization homology, and Ayala--Francis do not internally
re-derive HA~4.8 (they import it as a black box for the tensor
product of $E_{n}$-algebras, but do not prove it).
\end{proof}

\begin{remark}[Ginot's ext-verification via factorization algebras]
\label{v4-arith:tate:rem:ginot-ext-verification}
A further independent verification comes from Ginot (arXiv:1309.6044
\emph{Notes on factorization algebras, factorization homology, and
applications}): the restricted tensor product can be computed as the
factorization algebra of a single $E_{n}$-algebra on the disjoint
union of little-disks, with insertion of a distinguished spherical
vector at each added disk. The resulting factorization algebra agrees
with both of the previous two constructions. Three disjoint routes,
one answer; the HZ-IV triangulation passes.
\end{remark}

\section{Remarks on Operadic Coherence}
\label{v4-arith:tate:coherence-remarks}

\subsection{Why $E_{n}$-coherence is required}

We used coherence of $E_{n}$ (HA~3.3.1.1) in two places:

\begin{enumerate}
\item \Cref{v4-arith:tate:thm:alg-O-presentable}(1): presentability of
$\mathrm{Alg}_{\mathcal{O}}(\mathcal{E})$.
\item \Cref{v4-arith:tate:prop:filtered-colim-operadic}: filtered
colimits in $\mathrm{Alg}_{\mathcal{O}}(\mathcal{E})$ are computed
underlyingly.
\end{enumerate}

Both of these require the coherence hypothesis on
$\mathcal{O}$, which amounts to the statement that $\mathcal{O}$ has
enough operations to express its own composition; $E_{n}$ satisfies
this by HA~5.1.0.7.

\subsection{Alternative Unital Structures}

We chose the spherical-unit-insertion definition of
$\iota_{S\subseteq S'}$ via central units. An alternative
unital structure is:

\begin{remark}[non-spherical insertion]
\label{v4-arith:tate:rem:non-spherical-insertion}
Instead of inserting $u_{v}$, one could insert any
$\mathcal{O}$-central object $\tilde{u}_{v}\in\mathcal{C}_{v}$.
The resulting restricted tensor product is canonically equivalent
to the one with spherical units when $\tilde{u}_{v}$ is equivalent
to $u_{v}$ as $\mathcal{O}$-central objects. Different choices can
give different restricted tensor products; the natural
spherical-unit choice is distinguished by the unit property.
\end{remark}

\subsection{Iterated Restricted Tensor Products}

\begin{proposition}[iterated restricted tensor products]
\label{v4-arith:tate:prop:iterated-restricted}
\ClaimStatusProvedHere. Let $V=V_{1}\sqcup V_{2}$ be a partition,
with spherical data restricting to each $V_{i}$. Then
\[
\bigotimes_{v\in V}^{\prime}(\mathcal{C}_{v},u_{v})\;\simeq\;\Bigl(\bigotimes_{v\in V_{1}}^{\prime}(\mathcal{C}_{v},u_{v})\Bigr)\otimes\Bigl(\bigotimes_{v\in V_{2}}^{\prime}(\mathcal{C}_{v},u_{v})\Bigr)
\]
as $\mathcal{O}$-monoidal $\infty$-categories.
\end{proposition}

\begin{proof}
Both sides are filtered colimits over
$\mathrm{Fin}^{\mathrm{sph}}(V)=\mathrm{Fin}^{\mathrm{sph}}(V_{1})\times\mathrm{Fin}^{\mathrm{sph}}(V_{2})$
(cofinal product of filtered posets). The RHS distributes the
product colimit via HA~5.5.2.9 (finite products of filtered
colimits). LHS computes the colimit directly. The two agree by
HTT~4.2.3.10 (colimits of tensor products in symmetric monoidal
$\infty$-categories).
\end{proof}

\section{Completeness of the Closure of F2}
\label{v4-arith:tate:completeness-F2}

\begin{theorem}[closure of F2]
\label{v4-arith:tate:thm:F2-closure}
\ClaimStatusProvedHere. F2
(\Cref{v4-arith:conj:F2-perfectoid-E2}) closes: there is a
canonical $E_{2}$-monoidal structure on
$\mathrm{Hk}^{\mathrm{Iw}}_{GL_{2},\mathrm{glob}}$ such that the
global Iwahori Hecke category equals the restricted tensor product
of local Iwahori Hecke categories over $\overline{C}$ with spherical
units $|0\rangle_{v}$ at every place. This structure is unique up
to contractible choice.
\end{theorem}

\begin{proof}
Combine (F2.a), (F2.b) from
\Cref{v4-arith:thm:F2-reduction} with
\Cref{v4-arith:tate:cor:F2-closed}. (F2.a)+(F2.b) supply the
local $E_{2}$-structures and their spherical units;
\Cref{v4-arith:tate:cor:F2-closed} applies
\Cref{v4-arith:tate:thm:F2c-main} to produce the global
$E_{2}$-structure, and \Cref{v4-arith:tate:prop:uniqueness}
produces the uniqueness.
\end{proof}

\begin{remark}[status update]
\label{v4-arith:tate:rem:F2-status-final}
Prior to this chapter F2 was conjectural; after
\Cref{v4-arith:tate:thm:F2-closure} F2 is
a theorem. The status update table
(\Cref{v4-arith:cl:table:status}) is updated: F2 moves from
``Reduction to $\infty$-Tate lemma ProvedHere'' to simply
``ProvedHere''. This unlocks downstream statements depending on
F2, including F3's Dunn-additivity reduction
(\Cref{v4-arith:thm:F3-reduction}) and the global
$E_{2}$-structure used in \Cref{v4-arith:cl:rem:F3-inertia} for
the arithmetic inertia circle.
\end{remark}

\begin{corollary}[F3 reduces to F1 plus arithmetic Ben-Zvi--Nadler]
\label{v4-arith:tate:cor:F3-reduces-with-F2-closed}
\ClaimStatusProvedHere. With F2 now closed by
\Cref{v4-arith:tate:thm:F2-closure}, the
closure of F3 (\Cref{v4-arith:thm:F3-reduction}) reduces to the
conjunction of F1 and (F3.BZN$^{\mathrm{Ar}}$) only; the F2
hypothesis is discharged.
\end{corollary}

\begin{proof}
\Cref{v4-arith:thm:F3-reduction} assumed F1 plus F2 plus
(F3.BZN$^{\mathrm{Ar}}$). By
\Cref{v4-arith:tate:thm:F2-closure}, F2 is now a theorem;
it need no longer be assumed. Hence the conjunction F1 +
(F3.BZN$^{\mathrm{Ar}}$) suffices for F3.
\end{proof}

\section{Summary}
\label{v4-arith:tate:summary}

\begin{enumerate}
\item The $\infty$-categorical Tate restricted tensor product of
$\mathcal{O}$-monoidal stable presentable $\infty$-categories with
distinguished spherical units inherits a canonical
$\mathcal{O}$-monoidal structure
(\Cref{v4-arith:tate:thm:F2c-main}); the proof is a direct
application of Lurie~HA~4.8 + 5.1 and the centrality of the unit
(HA~5.1.2.26).
\item The statement is place-agnostic: it applies to any indexing
set $V$ with cofinite-many spherical units. The arithmetic
application $V=|\overline{C}|$, $\mathcal{C}_{v}=\mathrm{Hk}^{\mathrm{Iw}}_{GL_{2},v}$,
$u_{v}=|0\rangle_{v}$ closes F2
(\Cref{v4-arith:tate:cor:F2-closed}).
\item The structure is unique up to contractible choice
(\Cref{v4-arith:tate:prop:uniqueness}); functorial in the family
of $\mathcal{O}$-categories
(\Cref{v4-arith:tate:prop:functoriality-family}); and
base-change-covariant along symmetric monoidal colimit-preserving
functors (\Cref{v4-arith:tate:prop:base-change}).
\item Dunn additivity
(\Cref{v4-arith:tate:thm:dunn-additivity}) commutes with the
restricted tensor product
(\Cref{v4-arith:tate:prop:dunn-compat});
iterated constructions (e.g.~$E_{1}$-algebras in $E_{1}$-categories
giving an $E_{2}$-restricted-tensor) agree with direct
$E_{2}$-application.
\item Independent verification via Ayala--Francis factorization
homology (\Cref{v4-arith:tate:thm:verification-via-factorization-homology})
and Ginot's factorization-algebra account
(\Cref{v4-arith:tate:rem:ginot-ext-verification}) confirm the
$\mathcal{O}$-monoidal structure through a second, disjoint route;
the HZ-IV source-disjointness requirement is met.
\item F2 closes unconditionally
(\Cref{v4-arith:tate:thm:F2-closure}). F3's reduction simplifies:
F1 + (F3.BZN$^{\mathrm{Ar}}$) suffice
(\Cref{v4-arith:tate:cor:F3-reduces-with-F2-closed}).
\end{enumerate}

The cost of F2 was a $\sim$25-page application of the
$\infty$-operadic infrastructure established by Lurie HA
(especially~4.8 and~5.1). No arithmetic-specific input was needed:
F2.c is a pure $\infty$-categorical statement, and the arithmetic
is encoded only in the choice of the place-set $V$ and the family
$\{\mathrm{Hk}^{\mathrm{Iw}}_{GL_{2},v}\}$. The arithmetic branch
thereby discharges its ``single open $\infty$-categorical lemma''
and moves F2 from conjecture to theorem.

The next step is F3, which requires the arithmetic
Ben-Zvi--Nadler identification of Hochschild trace with functions
on the arithmetic free loop space. That is a genuinely arithmetic
obstruction; it is not absorbed into the $\infty$-categorical
infrastructure of Lurie HA, and it remains open
(\Cref{v4-arith:thm:F3-reduction}).
